% !TeX root = main.tex
\changecolours{ScoutGreen}
\section[CRUD com Python (PyMongo)]{Programando com o Driver Oficial Python (\texttt{PyMongo})}

% ------------------------------------------------------------------------------
\begin{frame}[fragile]{Conectando o Python ao MongoDB e Inserindo Documentos}
  \vspace{0.05cm}
  \begin{columns}[t]
    \begin{column}{0.48\textwidth}
      \alert{1. Instalação e Teste de Conexão (\texttt{ping})}\par
      \begin{lstlisting}[language=Python]
# No terminal: pip install pymongo
from pymongo import MongoClient
from pprint import pprint

# Conectar ao servidor local na porta 27017:
uri = "mongodb://localhost:27017/"
client = MongoClient(uri)

# Confirmar conexao com o comando ping:
client.admin.command("ping")
print("Conectado ao MongoDB com sucesso!")

# Selecionar banco 'mongodb_books' e colecao 'books'
books = client.mongodb_books.books
      \end{lstlisting}
    \end{column}
    \begin{column}{0.48\textwidth}
      \alert{2. Inserindo com \texttt{insert\_one} e \texttt{insert\_many}}\par
      \begin{lstlisting}[language=Python]
# Inserir 1 documento (dicionario Python):
book = {
    "isbn": "301",
    "name": "Python and MongoDB",
    "meta": {"version": "MongoDB 7.0"},
    "price": 60
}
res = books.insert_one(book)
print("ID gerado:", res.inserted_id)

# Inserir multiplos documentos em lote:
books.insert_many([
    {"isbn": "302", "name": "Python and MongoDB",
     "meta": {"version": "MongoDB 7.0"}, "price": 60},
    {"isbn": "303", "name": "Advanced MongoDB Techniques",
     "price": 70}
])
      \end{lstlisting}
    \end{column}
  \end{columns}
\end{frame}

% ------------------------------------------------------------------------------
\begin{frame}[fragile]{Consultando Documentos no Python: \texttt{find}, \textit{Dot Notation}, \texttt{\$gt} e \texttt{\$or}}
  \vspace{0.05cm}
  \begin{columns}[t]
    \begin{column}{0.48\textwidth}
      \alert{1. Busca Simples, \textit{Dot Notation} e \texttt{\$regex}}\par
      \begin{lstlisting}[language=Python]
# 1. Buscar pelo nome do livro:
for doc in books.find({"name": "Python and MongoDB"}):
    pprint(doc)

# 2. Buscar dentro do subdocumento 'meta' usando
# Dot Notation e Expressao Regular:
query_regex = {
    "meta.version": {
        "$regex": r".*?g.*?7\.0$",
        "$options": "i"
    }
}
for doc in books.find(query_regex):
    pprint(doc)
      \end{lstlisting}
    \end{column}
    \begin{column}{0.48\textwidth}
      \alert{2. Comparação (\texttt{\$gt}) e Operadores Lógicos}\par
      \begin{lstlisting}[language=Python]
# 3. Livros com preco maior que 50 ($gt):
for doc in books.find({"price": {"$gt": 50}}):
    pprint(doc)

# 4. Condicao logica AND (multiplas chaves):
q_and = {"name": "Advanced MongoDB Techniques",
         "price": 70}
pprint(books.find_one(q_and))

# 5. Condicao logica OR ($or):
q_or = {
    "$or": [
        {"name": "Advanced MongoDB Techniques"},
        {"isbn": "301"}
    ]
}
for doc in books.find(q_or):
    pprint(doc)
      \end{lstlisting}
    \end{column}
  \end{columns}
\end{frame}

% ------------------------------------------------------------------------------
\begin{frame}[fragile]{Atualizando (\texttt{update\_one}) e Removendo (\texttt{delete\_one}) no Python}
  \vspace{0.05cm}
  \begin{columns}[t]
    \begin{column}{0.48\textwidth}
      \alert{1. Atualizando com \texttt{update\_one} e \texttt{\$set}}\par
      \begin{lstlisting}[language=Python]
# Atualizar o preco do livro "Advanced MongoDB
# Techniques" para 75:
update_result = books.update_one(
    {"name": "Advanced MongoDB Techniques"},
    {"$set": {"price": 75}}
)

print("Encontrados:", update_result.matched_count)
print("Modificados:", update_result.modified_count)

# Buscar o documento atualizado para confirmar:
updated_doc = books.find_one(
    {"name": "Advanced MongoDB Techniques"}
)
pprint(updated_doc)
      \end{lstlisting}
    \end{column}
    \begin{column}{0.48\textwidth}
      \alert{2. Removendo com \texttt{delete\_one}}\par
      \begin{lstlisting}[language=Python]
# Remover o livro de ISBN '303':
isbn_to_delete = "303"
delete_result = books.delete_one(
    {"isbn": isbn_to_delete}
)
print("Removidos:", delete_result.deleted_count)

# Listando colecoes do banco via Python:
db = client.mongodb_books
print("Colecoes:", db.list_collection_names())
      \end{lstlisting}
      \scriptsize
      $\bullet$ \textbf{Sincronia com o Compass:} Clique em \textbf{Refresh} na aba \textbf{Documents} do Compass para ver as alterações!
    \end{column}
  \end{columns}
\end{frame}
